\documentclass[11pt]{article}
\usepackage[margin=1in]{geometry}
\usepackage{amsmath, amssymb}
\usepackage{graphicx}
\usepackage{booktabs}
\usepackage{setspace}
\usepackage{tcolorbox}
\usepackage{hyperref}

\setstretch{1.25}
\setlength{\parindent}{0pt}

\newtcolorbox{examplebox}{
  colback=gray!6,
  colframe=black!60,
  title=Worked Numerical Example,
  fonttitle=\bfseries
}

\newtcolorbox{keybox}{
  colback=blue!4,
  colframe=blue!60,
  title=Key Concept,
  fonttitle=\bfseries
}


\begin{document}




    {\LARGE \textbf{Housing Affordability and Real Estate Market in Canada}}\\[0.5em]
    {\large ECON 204 -- Contemporary Economic Issues}\\[0.5em]
 Winter 2026\\


\hrule
\vspace{1em}

%=====================================================
\section{Introduction}

Housing affordability has become one of the most pressing economic challenges in Canada. Unlike many macroeconomic issues that affect households indirectly, housing costs directly shape daily living standards, long-term financial security, and wealth accumulation. Housing is the largest single expenditure for most households and the dominant asset on household balance sheets. As a result, changes in housing affordability affect inequality, intergenerational mobility, and financial stability.

Since the mid-2000s,and especially after the COVID-19 pandemic,housing prices and rents in Canada have risen faster than incomes in many regions. Although higher interest rates after 2022 slowed price growth, affordability has remained strained because borrowing costs increased sharply and rental markets tightened. This chapter develops a detailed economic framework to explain these developments, drawing on Canadian institutions, quantitative measures, and international comparisons.

%=====================================================
\section{What Is Housing Affordability?}

\begin{keybox}
Housing affordability refers to the ability of households to obtain adequate housing without excessive financial strain, given their income, wealth, and access to credit.
\end{keybox}

Affordability is a \emph{relationship}, not a price level. High house prices do not automatically imply unaffordability if incomes are high or financing costs are low. Conversely, moderate prices can be unaffordable if incomes are low or credit conditions are tight.

\subsection*{Price-to-income measures}

A widely used indicator is the price-to-income ratio:
\[
PTI = \frac{\text{House Price}}{\text{Annual Household Income}}.
\]

\begin{examplebox}
Suppose the median house price in a Canadian city is \$900{,}000 and the median household income is \$110{,}000. The price-to-income ratio is:
\[
PTI = \frac{900{,}000}{110{,}000} \approx 8.18.
\]
This implies that the price of a typical home is more than eight times annual household income, a level widely regarded as severely unaffordable.
\end{examplebox}

Historically, PTI ratios around 3–4 were common in many Canadian cities. Ratios above 7–8 reflect a breakdown in the traditional link between incomes and housing prices.

\subsection*{Payment-based affordability}

Many households experience affordability through monthly cash flow rather than prices. Payment-based measures focus on the shelter-cost-to-income ratio:
\[
\text{Shelter Burden} = \frac{\text{Monthly Housing Costs}}{\text{Monthly Income}}.
\]

\begin{examplebox}
A household earns \$6{,}000 per month after tax. If mortgage payments, property taxes, and utilities sum to \$2{,}700 per month, the shelter burden is:
\[
\frac{2{,}700}{6{,}000} = 45\%.
\]
This household is housing-cost burdened even if the home value is high.
\end{examplebox}

This measure is especially important in Canada, where mortgage payments can change significantly at renewal.

%=====================================================
\section{Housing as a Consumption Good and an Asset}

Housing is unusual because it plays two distinct economic roles at the same time.

First, housing is a consumption good. Households consume housing services every day in the form of shelter, space, safety, and location. From this perspective, housing is similar to other goods and services: people pay for it using rent or implicit rent (if they own), and demand depends on income, household size, and preferences. If housing were only a consumption good, we would expect housing prices to be closely tied to fundamentals such as rents and household incomes, because these determine how much households can and are willing to pay for shelter services.

Second, housing is also a store of wealth and an investment asset. When households buy a home, they are not only purchasing shelter for today; they are acquiring an asset that may appreciate in value over time. For many Canadian households, housing is the largest component of net worth. Expectations about future house prices therefore influence current housing demand in the same way that expectations about future stock prices influence stock market demand.

This dual role,consumption plus investment,means that housing demand depends not only on current ability to pay for shelter, but also on expected future returns. When households expect house prices to rise, owning a home becomes attractive even if current prices seem high relative to rents or incomes. Buyers may be willing to stretch financially today because they expect capital gains tomorrow.

As a result, housing prices can diverge from traditional fundamentals such as rents or incomes. For example, rents may rise slowly because rental supply is relatively stable or regulated, while house prices rise rapidly because buyers anticipate future price appreciation. Similarly, incomes may grow modestly while house prices surge if low interest rates and optimistic expectations make housing appear to be a good investment.

A simple way to formalize this idea is through the user cost of housing framework. The cost of owning a home in a given year depends on interest costs, taxes, and maintenance, but it is reduced by expected house price appreciation. When expected appreciation is high, the effective cost of ownership falls, encouraging higher demand and higher prices even if rents and incomes do not change.

A concrete example helps illustrate the point. Suppose two cities have identical rents and incomes. In City A, households expect house prices to be stable. In City B, households expect prices to rise by 5 percent per year. Even though the shelter services are the same, demand for ownership will be much higher in City B because buyers expect capital gains. As a result, house prices in City B will be higher, and the price-to-rent ratio will be larger, even though rental fundamentals are identical.

This feature helps explain many real-world housing booms. During periods of low interest rates or strong optimism,such as Canada in the late 2010s and early 2020s,house prices increased much faster than rents or incomes. The divergence does not require irrational behavior; it can arise from rational expectations about future prices combined with the asset nature of housing.

In short, housing markets can move away from rents and incomes because housing is not just a place to live, but also a financial asset whose value depends on expectations about the future. This dual role makes housing markets especially sensitive to interest rates, expectations, and financial conditions, and helps explain why housing affordability can deteriorate rapidly even when current economic fundamentals appear stable.




%=====================================================
\section{Demand-Side Drivers of Housing Affordability in Canada}

Demand-side drivers of housing affordability refer to the factors that increase households’ willingness and ability to pay for housing at any given price. In Canada, these forces have been unusually strong and persistent, helping to explain why housing costs have risen faster than incomes in many regions even during periods of elevated construction activity. The effect of demand-side pressures on affordability depends critically on the responsiveness of housing supply. When supply adjusts slowly, as is often the case in Canadian cities, increases in demand are reflected primarily in higher prices and rents rather than in expanded housing availability.

\subsection*{Population growth and household formation}

One of the most important demand-side drivers of housing affordability in Canada is rapid population growth. Over the past decade, Canada has experienced among the fastest population growth rates in the advanced world, driven by immigration, temporary residents, international students, and natural population increase. Population growth raises housing demand immediately, as new households require housing services upon arrival.

Affordability deteriorates when the rate of household formation exceeds the rate at which housing units are completed. This effect is particularly pronounced in large metropolitan areas, where population inflows are concentrated. Even when population growth supports long-run economic growth, the short-run impact is often intensified competition for existing housing stock, leading to higher prices and rents.

\subsection*{Income growth and income dispersion}

Housing affordability is shaped not only by average income growth, but also by the distribution of income. In Canada, higher-income households have experienced stronger income growth than median households, particularly among dual-earner families and workers in high-paying sectors such as finance, technology, and professional services. These households have greater bidding power in housing markets, especially in desirable locations.

As a result, housing prices often reflect the purchasing power of households in the upper part of the income distribution rather than that of the median household. This can lead to widespread affordability challenges even when aggregate income indicators appear strong. Middle-income households may find themselves priced out despite stable employment and rising nominal incomes.

\subsection*{Credit conditions and interest rates}

Mortgage credit conditions play a central role in shaping housing demand. Housing affordability depends not only on income, but also on how much households can borrow. Lower interest rates increase the size of the mortgage that can be serviced at a given monthly payment, raising effective housing demand.

In Canada, extended periods of low interest rates,particularly following the Global Financial Crisis and during the COVID-19 pandemic,substantially increased borrowing capacity. This mechanical effect allowed households to bid up housing prices even in the absence of strong income growth. Conversely, increases in interest rates reduce borrowing capacity but may worsen affordability in the short run by raising mortgage payments at renewal and pushing more households into the rental market, thereby increasing rental demand and rents.

\subsection*{Expectations of future house prices}

Expectations about future house prices are an important determinant of housing demand. When households expect house prices to rise, owning a home becomes more attractive both as a source of shelter and as an investment. Expected price appreciation reduces the perceived long-run cost of ownership and can encourage households to enter the market sooner or to accept higher levels of financial risk.

In Canada, sustained periods of price growth have reinforced expectations of continued appreciation, particularly in major urban markets. This expectation channel can generate self-reinforcing demand, as rising prices attract additional buyers, which in turn pushes prices higher. Such dynamics can arise even among forward-looking and risk-aware households.

\subsection*{Demographic and life-cycle factors}

Demographic structure influences housing demand through life-cycle patterns of household formation and tenure choice. Large cohorts entering prime home-buying ages increase demand for ownership housing, while longer life expectancy and preferences for aging in place reduce turnover in the existing housing stock. In Canada, these factors have contributed to sustained demand pressure in ownership markets.

At the same time, affordability constraints delay homeownership for younger households, increasing demand in the rental sector. This shift places upward pressure on rents and tightens rental markets, worsening affordability even for households not actively attempting to purchase homes.

\subsection*{Urban concentration and location preferences}

Housing demand in Canada is increasingly concentrated in large urban centres where employment opportunities, amenities, and public services are clustered. Even when national housing supply expands, affordability can deteriorate locally if demand growth is spatially concentrated and supply responses are slow in high-demand areas.

Urban concentration also generates spillover effects, as households priced out of central cities relocate to surrounding regions, transmitting demand pressures across metropolitan areas and contributing to broader regional affordability challenges.






%=====================================================
\section{Supply Constraints and Price Dynamics}

Supply constraints play a central role in explaining why housing affordability deteriorates sharply when demand increases. In housing markets, prices are determined by the interaction between demand and the ability of housing supply to expand. When supply responds slowly or imperfectly, increases in demand are absorbed primarily through higher prices and rents rather than through the construction of additional housing units. This mechanism is fundamental to understanding housing price dynamics in Canada and in other advanced economies.

Housing supply is inherently slower to adjust than supply in many other markets. Housing is durable, land-intensive, and heavily regulated. The production of new housing units requires land assembly, zoning approval, financing, construction labour, and supporting infrastructure, all of which take time. Even in markets with strong incentives to build, supply responses typically occur over years rather than months. As a result, short-run housing supply is relatively fixed, and prices must rise to clear the market when demand increases.

A key concept for understanding supply-side dynamics is \textbf{housing supply elasticity}, which measures how responsive housing construction is to changes in prices. In markets with elastic supply, higher prices lead to substantial increases in construction. In markets with inelastic supply, higher prices generate little additional construction, and adjustment occurs almost entirely through prices. Empirical evidence indicates that housing supply in many Canadian metropolitan areas,particularly large cities,is relatively inelastic.

Land-use regulation is one of the most important sources of supply inelasticity. Zoning rules that limit density, restrict building heights, or reserve large areas for single-family housing constrain the number of housing units that can be built even when demand is strong. Lengthy approval processes and uncertainty surrounding permitting further discourage rapid supply expansion. When developers face delays or unpredictable outcomes, they are less able to respond quickly to rising prices, reinforcing upward pressure on housing costs.

Geographic constraints also play an important role in shaping supply responsiveness. Physical features such as coastlines, mountains, or protected land reduce the amount of developable land in some cities. In these environments, expanding housing supply requires either costly redevelopment or increased density, both of which often face political and regulatory resistance. Consequently, demand shocks such as population growth or declines in interest rates are more likely to translate into higher prices than increased housing quantities.

Construction capacity imposes an additional constraint on supply. Even when zoning permits development, shortages of skilled labour, rising material costs, and limited infrastructure capacity can slow the pace of construction. These constraints become particularly binding during periods of rapid demand growth, when the construction sector is unable to scale up output sufficiently to offset rising housing demand.

The interaction between demand shocks and supply constraints helps explain why housing prices can rise sharply even when construction activity appears robust. High levels of housing starts do not necessarily indicate adequate supply if completions lag behind demand or if population growth continues to outpace housing delivery. In such cases, prices continue to rise because the effective supply available to households remains insufficient.

Supply constraints also contribute to the persistence of high housing prices. When supply responds slowly, prices do not quickly revert to previous levels after a demand shock dissipates. Instead, elevated prices become embedded in the market, particularly if households and investors expect supply constraints to persist. This persistence helps explain why housing affordability often fails to improve substantially even after housing booms cool or monetary policy tightens.

Finally, supply constraints shape the distributional consequences of housing market pressures. When new housing is scarce, access is rationed through prices. Higher-income households are better positioned to absorb price increases, while lower- and middle-income households are pushed into smaller dwellings, longer commutes, or the rental market. In this way, supply constraints not only raise average housing costs but also exacerbate inequality and spatial segregation within cities.

In summary, supply constraints are a critical driver of housing price dynamics. When housing supply is slow to respond due to regulatory, geographic, and capacity limitations, increases in demand lead primarily to higher prices and rents rather than increased housing availability. Understanding these supply-side mechanisms is essential for explaining persistent affordability challenges and for evaluating policies aimed at improving housing affordability.


%=====================================================
\section{Canadian Housing Market Case Studies}

Examining housing affordability through Canadian case studies is essential because housing markets in Canada are highly regional. National averages often conceal large differences in prices, rents, supply conditions, and institutional constraints. By comparing cities and regions, it becomes clear how common economic forces,such as population growth, interest rates, and income dynamics,interact differently depending on local supply elasticity, geography, and policy environments. These case studies also demonstrate that there is no single Canadian housing market, but rather a collection of distinct local markets responding to shared national and global pressures.

\subsection*{Vancouver: Inelastic Supply and Price Capitalization}

Vancouver represents a textbook example of a housing market characterized by strong demand and highly inelastic supply. Demand is supported by high incomes, desirable amenities, and sustained population growth. At the same time, supply is constrained by a combination of geographic limitations, including mountains and ocean, limited developable land, and restrictive land-use regulations that limit density across much of the metropolitan area.

Because supply responds slowly, increases in demand are reflected primarily in higher prices rather than in expanded housing stock. Even relatively small demand shocks,such as lower interest rates or population inflows,can therefore generate large price increases. This dynamic explains why Vancouver has persistently high price-to-income ratios and why affordability has remained strained even during periods of strong construction activity.

Vancouver also illustrates how housing prices can become increasingly disconnected from local incomes. Prices are influenced not only by local shelter demand but also by expectations of future price appreciation, which increases the attractiveness of housing as an investment asset. As a result, price growth has at times outpaced both rent growth and income growth, consistent with the implications of low supply elasticity.

\subsection*{Toronto and the Greater Golden Horseshoe: Growth Outpacing the Pipeline}

Toronto’s housing affordability challenges are closely linked to rapid population growth combined with long development timelines. The Greater Golden Horseshoe has absorbed a substantial share of Canada’s population increase, driven by immigration and employment opportunities. This has created sustained demand across both ownership and rental markets.

Unlike Vancouver, Toronto does not face severe geographic constraints. However, housing supply is limited by zoning restrictions, infrastructure coordination challenges, and lengthy approval processes. High-rise construction is common, yet large projects often take many years to move from approval to completion. As a result, housing starts may appear strong while completed units lag behind demand, maintaining upward pressure on prices and rents.

Toronto’s experience highlights an important insight: even when supply is expanding, affordability may continue to deteriorate in the short to medium run if demand growth exceeds the rate at which housing can be delivered.

\subsection*{Montreal: From Relative Affordability to Tightening Markets}

Montreal has historically been viewed as a relatively affordable large city, particularly in the rental market. This reflected a large rental housing stock, slower population growth, and institutional features such as stronger tenant protections and more flexible density in some neighbourhoods.

In recent years, however, affordability pressures have intensified. Population growth accelerated, prolonged low interest rates increased housing demand, and supply did not adjust quickly enough. Although higher interest rates after 2022 moderated price growth, affordability challenges persisted as rental demand strengthened and mortgage payments increased at renewal.

Montreal’s experience demonstrates that relative affordability advantages are not permanent. Cities that appear affordable at one point in time can experience rapid deterioration when demand conditions shift and supply responses are slow.

\subsection*{Prairie Cities: Elastic Supply with Emerging Pressures}

Prairie cities such as Calgary and Edmonton have traditionally offered better housing affordability due to more elastic housing supply. Abundant land availability, lower density restrictions, and faster permitting processes allowed construction to respond more quickly to changes in demand. As a result, price-to-income ratios in these cities remained lower than in major metropolitan areas for many years.

Recently, however, strong interprovincial migration and population growth have increased housing demand sharply. Although supply has responded more than in constrained markets, it has not been sufficient to prevent rising prices and rents. This illustrates that even in relatively elastic markets, affordability can deteriorate when demand growth accelerates rapidly.

\subsection*{Atlantic Canada: Small Markets and Large Shocks}

Housing markets in Atlantic Canada, including Halifax and surrounding regions, experienced sharp affordability declines during the pandemic period. These markets are smaller and historically characterized by modest price growth. When remote work, migration, and demographic shifts increased housing demand, supply was slow to adjust due to limited construction capacity and smaller development pipelines.

Because initial price levels were lower, percentage price increases appeared especially large. This case illustrates how smaller markets can experience outsized affordability shocks when demand increases rapidly relative to existing housing stock.

\subsection*{Lessons from the Canadian Experience}

Taken together, these Canadian case studies illustrate several fundamental principles of housing economics. Housing supply elasticity plays a central role in determining price responses to demand shocks. Affordability dynamics are local, even when driven by national or global forces. Demand pressures can spread across regions as households relocate from more expensive cities to more affordable ones. Finally, affordability challenges can emerge quickly and persist when supply responses are slow.



%=====================================================
\section{Monetary Policy and Mortgage Renewal}

Monetary policy affects housing affordability through several channels, but the mortgage renewal channel is particularly important in Canada. This channel helps explain why housing affordability can worsen even when house prices stop rising or begin to fall. The key reason is that Canadian mortgages expose households to interest rate risk at renewal, rather than locking in borrowing costs for the entire life of the loan.

In Canada, most mortgages have fixed interest rates for relatively short terms, commonly five years, while amortization periods are much longer, often 25 or 30 years. As a result, households must periodically renew their mortgages at prevailing market interest rates. When the central bank raises policy interest rates, this does not immediately affect all existing mortgage payments. Instead, the impact arrives gradually as mortgages come up for renewal.

Monetary policy tightening raises the policy rate, which feeds into higher short-term interest rates and, through financial markets, into higher mortgage rates. When households renew their mortgages at these higher rates, their monthly payments can increase substantially, even if the outstanding loan balance has declined only modestly. This payment increase directly raises housing costs and worsens affordability through a cash-flow channel.

This mechanism is distinct from the price channel of monetary policy. Higher interest rates tend to reduce housing demand and slow house price growth by lowering borrowing capacity for new buyers. However, for existing homeowners with mortgages, higher interest rates increase required payments. As a result, housing prices and housing affordability can move in opposite directions. House prices may stabilize or decline, while affordability worsens because monthly payments rise.

Mortgage renewal effects are especially powerful because housing costs represent a large share of household income. Even moderate increases in interest rates can translate into hundreds or thousands of dollars per month in additional payments. For many households, wages do not adjust quickly enough to offset these increases, leading to higher shelter-cost-to-income ratios and increased financial stress.

The mortgage renewal channel also has important distributional consequences. Households that purchased homes recently or took on large mortgages during periods of low interest rates are more exposed to renewal risk. These households tend to be younger and more highly leveraged. In contrast, long-time homeowners with small remaining mortgage balances or no mortgage at all are largely insulated from rising interest rates. As a result, monetary tightening can redistribute financial stress across households and exacerbate intergenerational inequality.

Monetary policy affects renters indirectly through the mortgage renewal channel as well. Higher interest rates raise financing costs for landlords, and in tight rental markets some of these higher costs may be passed on to renters through higher rents. In addition, higher interest rates can delay new housing construction, constraining future rental supply and reinforcing upward pressure on rents.

An important implication of the mortgage renewal channel is that the effects of monetary policy on housing unfold with long and uneven lags. Even after interest rates stop rising, affordability pressures may continue to intensify as more households renew their mortgages at higher rates. This helps explain why housing-related financial stress can peak well after a tightening cycle begins.

In summary, monetary policy influences housing affordability in Canada not only by affecting house prices, but also by changing mortgage payments at renewal. Because Canadian mortgages reset periodically, interest rate increases translate into higher housing costs for existing homeowners over time. This renewal mechanism explains why housing affordability can deteriorate even when house prices are stable or falling, and it highlights the importance of payment-based measures when evaluating the housing effects of monetary policy.

%=====================================================
*{International Housing Affordability Crises}

The term \textbf{international housing affordability crisis} refers to the widespread and simultaneous deterioration of housing affordability across many advanced economies over the past two decades. What distinguishes the current episode from earlier housing booms is that affordability problems have emerged across a wide range of countries with very different institutions, housing systems, and policy frameworks. This pattern suggests that the crisis is driven by common economic forces rather than by isolated national policy failures.

At the core of the international housing affordability crisis is a persistent divergence between housing costs and household incomes. In many countries, house prices and rents have increased substantially faster than median incomes, pushing traditional affordability metrics such as price-to-income ratios and rent-to-income ratios to historically high levels. In major cities around the world, purchasing a home now requires many multiples of annual income, while renting absorbs an increasing share of household earnings.

\subsection*{A common global pattern}

Despite institutional differences, countries experiencing affordability crises share several key characteristics. First, housing demand is increasingly concentrated in major urban centres where employment opportunities, amenities, and public services are clustered. Second, housing supply in these cities responds slowly due to zoning restrictions, land-use regulation, infrastructure constraints, and construction bottlenecks. Third, long periods of favourable financial conditions,particularly low interest rates,have increased borrowing capacity and raised asset prices.

Because housing supply tends to be relatively inelastic in dense urban areas, increases in demand translate primarily into higher prices rather than increased quantities. As a result, affordability deteriorates even when overall economic conditions appear strong.

\subsection*{The United States: regional divergence}

The United States illustrates how supply responsiveness shapes affordability outcomes. In cities such as San Francisco, Los Angeles, and New York, housing supply is highly constrained by zoning rules, community opposition to development, and lengthy approval processes. In these markets, rising demand has led mainly to higher prices rather than increased construction, producing severe affordability pressures.

By contrast, cities with more flexible land-use regulations,such as Houston or many Midwestern cities,have experienced smaller increases in house prices. This contrast highlights a central lesson of the international housing affordability crisis: where housing supply is inelastic, demand shocks generate large price increases and rapid declines in affordability.

\subsection*{The United Kingdom: rental affordability and generational inequality}

In the United Kingdom, housing affordability challenges are particularly acute in the rental market. Rising house prices have made homeownership increasingly unattainable for younger households, pushing more people into renting. As rental demand has grown faster than rental supply, rents have risen sharply, especially in London and the South East.

An important distributional feature of the UK experience is that renters face higher effective inflation than homeowners. Homeowners with fixed-rate mortgages are partially insulated from rising housing costs, while renters experience immediate rent increases. This dynamic has intensified generational inequality, with younger cohorts devoting a much larger share of income to housing than older homeowners.

\subsection*{Australia and New Zealand: extreme price-to-income ratios}

Australia and New Zealand are frequently cited as among the most unaffordable housing markets in the world. In cities such as Sydney, Melbourne, and Auckland, price-to-income ratios have reached exceptionally high levels. These outcomes reflect a combination of strong population growth, urban concentration of economic activity, constrained housing supply, and financial conditions that have favoured housing as an investment asset.

These cases demonstrate that housing affordability crises can persist even in high-income countries with stable macroeconomic institutions and well-functioning financial systems.

\subsection*{Europe: different institutions, similar outcomes}

European countries vary widely in housing tenure systems, with some emphasizing rental housing and others ownership. Nevertheless, affordability problems have emerged across many European cities. In metropolitan areas such as Paris, Amsterdam, and Stockholm, limited housing supply combined with strong demand has pushed prices and rents far above income growth.

The convergence of affordability challenges across diverse European housing systems suggests that the crisis is not driven by a single institutional model, but by the interaction of global demand pressures with local supply constraints.

\subsection*{Why affordability has become a structural problem}

Although housing markets have always been cyclical, the international affordability problem is increasingly described as a structural crisis rather than a temporary fluctuation. Affordability has worsened persistently over time, younger and lower-income households are increasingly excluded from homeownership, and housing costs are reshaping labour mobility, family formation, and inequality.

In many countries, affordability has not returned to pre-boom levels even after housing markets cooled, indicating that the underlying forces driving high housing costs remain in place.

\subsection*{Implications for Canada}

Canada’s housing affordability challenges fit squarely within this international pattern. Like other advanced economies, Canada combines strong demand concentrated in large cities, constrained housing supply, and financial conditions that amplify price movements. Viewing Canada’s experience through an international lens helps clarify that the problem is not unique, but part of a broader transformation in housing markets across advanced economies.




%=====================================================
\section{Conclusion}

Housing affordability in Canada reflects the interaction of strong demand growth, constrained supply, and powerful financial channels. Prices, rents, and payments respond differently to economic shocks, meaning affordability can deteriorate even when some indicators improve. Understanding these mechanisms is essential for evaluating policy responses and for interpreting housing market developments in Canada and abroad.


\newpage
\section*{References}

Bank of Canada. 2025. “Mortgage Renewals and Payment Impacts.” \textit{Staff Analytical Note} 2025--21. Ottawa: Bank of Canada.

Canada Mortgage and Housing Corporation (CMHC). 2025. \textit{Housing Supply Gaps Report}. Ottawa: CMHC.

Canada Mortgage and Housing Corporation (CMHC). 2025. \textit{Housing Supply Report: Market Data and Research}. Ottawa: CMHC.

Organisation for Economic Co-operation and Development (OECD). 2025. \textit{OECD Economic Surveys: Canada 2025}. Paris: OECD Publishing.

Pomeroy, Steve, and Frank Clayton. 2025. “Zoning Constraints, Housing Supply, and Affordability in Canadian Cities.” \textit{Canadian Public Policy} 51 (1): 1–25.

Demographia. 2025. \textit{Demographia International Housing Affordability}. St. Louis: Demographia.

The Guardian. 2025. “Australian Households Spend a Growing Share of Income on Mortgages.” November 24.




\end{document}
